\chapter{Finite-Dimensional Operator Foundations}
\label{ch:finite-operator-foundations}

\section{The finite retained-space problem}

Finite-dimensional reduction replaces a large operator problem by its action on a
smaller declared subspace. Let $\mathcal H_{\mathrm{amb}}\cong\mathbb C^D$ be an ambient
Hilbert space carrying an operator $H$, and let
$\mathcal H_{\mathrm{ret}}\cong\mathbb C^M$, with $M\leq D$, be a retained coordinate
space. A rectangular matrix $V\in\mathbb C^{D\times M}$ maps retained-coordinate vectors
into the ambient space. Its columns may be selected eigenvectors, projected trial
states, reduced-basis vectors, or any other declared finite frame.

Two mathematical objects must be kept distinct. The range of $V$ is the retained
subspace inside the ambient Hilbert space. The ordered columns of $V$ are one coordinate
frame for that subspace. When the columns are orthonormal, the subspace is represented
in ambient coordinates by
\[
  P_V=VV^\dagger,
\]
whereas the matrix of $H$ in retained coordinates is
\[
  C_V(H)=V^\dagger HV.
\]
The first object answers which ambient vectors are retained. The second answers how the
ambient operator acts after it has been expressed in a selected retained basis.

The distinction matters because a retained subspace has no unique ordered basis. If
$G\in U(M)$, then $VG$ spans the same subspace as $V$. A sound reduction formalism must
therefore show that the ambient projector is unchanged, identify the covariance law for
the compressed matrix, and verify that reconstruction in the ambient space does not
depend on the retained-coordinate basis. Otherwise a change of coordinates can be
mistaken for a change of model.

The ten theorems in this chapter establish those finite identities without invoking
periodicity, crystal momentum, Fourier transformation, or spatial locality. They form a
representation-neutral foundation that can be reused for Bloch frames, trial subspaces,
finite quantum systems, and other matrix reductions. The later Bloch chapter adds the
extra structure carried by a momentum-indexed family of frames.

\section{Definitions and notation}

Let the ambient and retained Hilbert spaces be
\[
  \mathcal H_{\mathrm{amb}}\cong\mathbb C^D,
  \qquad
  \mathcal H_{\mathrm{ret}}\cong\mathbb C^M,
  \qquad M\leq D.
\]
The Lean declarations use arbitrary finite index types rather than natural-number
indices; $D$ and $M$ denote their cardinalities in the standard matrix notation used
here.  The modeling notation assumes $M\leq D$, but the Lean definitions and theorem
statements impose no separate cardinality inequality; exact frame orthonormality appears
as an explicit equation wherever it is needed.  The symbol $\dagger$ denotes conjugate
transpose, and $I_D$ and $I_M$ denote the identity matrices on the ambient and retained
spaces.

In Dirac notation, choose the coordinate basis
$\{|a\rangle_{\mathrm{ret}}\}_{a=1}^M$ of the retained space and let
$\{|v_a\rangle\}_{a=1}^M$ be the retained vectors in the ambient space. The frame matrix
is the linear embedding
\[
  V=\sum_{a=1}^M |v_a\rangle\,{}_{\mathrm{ret}}\!\langle a|,
  \qquad
  V^\dagger=\sum_{a=1}^M |a\rangle_{\mathrm{ret}}\langle v_a|.
\]
Thus $V$ is not one state when $M>1$; it is the ordered collection of retained kets,
viewed as a map from retained coordinates into the ambient Hilbert space. The identities
\[
  V^\dagger V=I_M,
  \qquad
  VV^\dagger=\sum_{a=1}^M |v_a\rangle\langle v_a|
\]
express, respectively, orthonormality of the retained kets and the associated ambient
projector. The compressed operator has matrix elements
\[
  \langle a|C_V(H)|b\rangle
  =\langle v_a|H|v_b\rangle.
\]
A right-unitary change $V\mapsto VG$ replaces the retained kets by
\[
  |v'_a\rangle=\sum_{b=1}^M |v_b\rangle G_{ba}.
\]
For $M=1$, the embedding reduces to a single normalized ket, $V=|v\rangle$, with
$P_V=|v\rangle\langle v|$ and $C_V(H)=\langle v|H|v\rangle$.

The variables used throughout this chapter are:

\begin{itemize}
  \item $V\in\mathbb C^{D\times M}$: a finite retained frame. Its columns are the
        retained vectors expressed in the ambient basis.
  \item $H\in\mathbb C^{D\times D}$: an ambient operator, typically interpreted as a
        finite Hamiltonian when $H^\dagger=H$.
  \item $G\in\mathbb C^{M\times M}$: a retained-coordinate unitary satisfying
        $G^\dagger G=GG^\dagger=I_M$.
  \item $V_G=VG$: the same frame after a right-unitary change of retained coordinates.
  \item $P_V=VV^\dagger\in\mathbb C^{D\times D}$: the ambient frame operator. When the
        columns of $V$ are orthonormal, $P_V$ is the orthogonal projector onto their
        span.
  \item $C_V(H)=V^\dagger HV\in\mathbb C^{M\times M}$: the compression of $H$ to
        retained coordinates.
  \item $R_V(H)=V C_V(H)V^\dagger\in\mathbb C^{D\times D}$: reconstruction of the
        compressed operator in the ambient coordinates.
\end{itemize}

\begin{definition}[Orthonormal finite frame]
The frame $V$ is orthonormal when
\[
  V^\dagger V=I_M.
\]
This is an exact hypothesis. It is not inferred from a floating-point tolerance.
\end{definition}

\begin{definition}[Hermitian ambient operator]
The operator $H$ is Hermitian when $H^\dagger=H$.
\end{definition}

\begin{remark}[Formal correspondence]
In \modulepath{ProjectKoios/BandGauge/Alignment/FiniteFrame.lean}, the definitions above
are named \decl{IsOrthonormalFrame}, \decl{frameProjector},
\decl{rightUnitaryFrame}, \decl{compressOperator}, and
\decl{reconstructCompressedOperator}. The prose proofs below explain the checked Lean
theorems but do not replace their proof terms.
\end{remark}

\section{Relation to the research manuscript's finite-dimensional foundation}

The working \texttt{ksdft2effmass} research manuscript uses the same representation-neutral
objects in its finite-dimensional-foundations and projection-and-retained-spaces
chapters. Table~\ref{tab:finite-manuscript-notation-bridge} records the notation bridge.

\begin{table}[ht]
\centering
\small
\begin{tabularx}{\textwidth}{@{}l X X@{}}
\toprule
This paper & Research manuscript & Common mathematical object \\
\midrule
$V$
  & $V$
  & $D\times M$ retained-frame embedding \\
$|v_a\rangle$
  & $|\phi_a\rangle$
  & column $a$ of the retained frame \\
$P_V=VV^\dagger$
  & $P$ or $\widehat P_{\mathcal S}$
  & ambient projector onto the retained subspace \\
$C_V(H)=V^\dagger HV$
  & $H_{\mathcal S}$
  & operator matrix in retained coordinates \\
$R_V(H)=VC_V(H)V^\dagger$
  & $PHP$ embedded in the parent space
  & reconstructed retained operator in ambient coordinates \\
$V_G=VG$
  & $V'=VG$
  & right-unitary change of retained frame \\
\bottomrule
\end{tabularx}
\caption{Representation-neutral notation bridge to the working research manuscript.}
\label{tab:finite-manuscript-notation-bridge}
\end{table}

The theorem correspondence is likewise direct. Declarations
\decl{BG-ALIGNMENT-THM-001}, \decl{BG-ALIGNMENT-THM-002},
\decl{BG-ALIGNMENT-THM-009}, and \decl{BG-ALIGNMENT-THM-010} justify the projector
properties used in the projection discussion. Declarations
\decl{BG-ALIGNMENT-THM-003} and \decl{BG-ALIGNMENT-THM-004} justify right-unitary frame
changes and projector invariance; \decl{BG-ALIGNMENT-THM-004} realizes the formula
presented prospectively as \texttt{PRF-05.01}. Declaration
\decl{BG-ALIGNMENT-THM-005} realizes compression covariance, presented prospectively as
\texttt{PRF-05.02}. The remaining results establish Hermiticity of compression,
projector-sandwich reconstruction, and reconstruction invariance.

This bridge is editorial, not authoritative. It records notation observed in the clean
\texttt{ksdft2effmass} worktree on branch
\decl{work/monograph-representation-carriers-and-derivation-graphs} at commit
\decl{e3d8a51b30e2072aa4bb116eb46389e342fc00fc}, principally in
\modulepath{chapters/finite-dimensional-foundations.tex} and
\modulepath{chapters/projection-and-retained-spaces.tex}. Those chapters supply the
source notation for this bridge. The Lean declarations identified here fix the exact
statements and proof status.

\section{Finite-frame theorem sequence}

\begin{theorem}[\texttt{BG-ALIGNMENT-THM-001}: Hermiticity of the frame operator]
\label{thm:bg-alignment-001}
For every finite complex frame $V\in\mathbb C^{D\times M}$,
\[
  P_V^\dagger=P_V.
\]
No orthonormality hypothesis is required.
\end{theorem}

\begin{proof}
By conjugate-transpose reversal and involution,
\[
  P_V^\dagger=(VV^\dagger)^\dagger
  =(V^\dagger)^\dagger V^\dagger
  =VV^\dagger=P_V.
\]
\end{proof}

Hermiticity is the first structural requirement for interpreting $P_V$ as an observable-
compatible subspace operator. Importantly, it holds before normalization is imposed.
This separates two logically different facts: every frame produces a Hermitian Gram-
type operator in the ambient space, whereas only an orthonormal frame makes that
operator an orthogonal projector. The distinction prevents approximate normalization
from being hidden inside a symmetry claim.

\begin{theorem}[\texttt{BG-ALIGNMENT-THM-002}: Idempotence of an orthonormal-frame projector]
\label{thm:bg-alignment-002}
If $V^\dagger V=I_M$, then
\[
  P_V^2=P_V.
\]
Together with Theorem~\ref{thm:bg-alignment-001}, this makes $P_V$ the orthogonal
projector onto the retained column space.
\end{theorem}

\begin{proof}
Associativity exposes the retained-space Gram matrix:
\[
  P_V^2=(VV^\dagger)(VV^\dagger)
  =V(V^\dagger V)V^\dagger
  =VI_MV^\dagger
  =VV^\dagger=P_V.
\]
\end{proof}

Together, Hermiticity and idempotence identify $P_V$ as the coordinate-free carrier of
the retained subspace. This is the object that should be compared when the scientific
question concerns which states are retained rather than how a basis for those states is
written. The theorem is also the point at which exact orthonormality becomes essential;
without it, $VV^\dagger$ need not be a projector.

\begin{theorem}[\texttt{BG-ALIGNMENT-THM-003}: Orthonormality after a retained unitary]
\label{thm:bg-alignment-003}
If $V^\dagger V=I_M$ and $G$ is unitary, then $V_G=VG$ is also an orthonormal frame:
\[
  V_G^\dagger V_G=I_M.
\]
\end{theorem}

\begin{proof}
Using conjugate-transpose reversal, the frame hypothesis, and $G^\dagger G=I_M$,
\[
  V_G^\dagger V_G
  =(VG)^\dagger(VG)
  =G^\dagger(V^\dagger V)G
  =G^\dagger I_MG
  =I_M.
\]
\end{proof}

Composite retained states admit internal unitary mixing, and numerical eigensolvers may
return different orthonormal bases for the same retained subspace. This theorem shows
that such a mixing remains inside the class of admissible orthonormal frames. Smoothness,
optimization, and physical selection enter only when a particular gauge is chosen.

\begin{theorem}[\texttt{BG-ALIGNMENT-THM-004}: Projector invariance under a retained unitary]
\label{thm:bg-alignment-004}
For every frame $V$ and every retained-space unitary $G$,
\[
  P_{V_G}=P_V.
\]
The conclusion does not require $V^\dagger V=I_M$.
\end{theorem}

\begin{proof}
The unitary and its adjoint cancel in the retained coordinates:
\[
  P_{V_G}=(VG)(VG)^\dagger
  =VGG^\dagger V^\dagger
  =VI_MV^\dagger
  =P_V.
\]
\end{proof}

The projector is unchanged even though the retained basis has changed. Consequently,
$V$ and $VG$ describe the same ambient subspace, not two competing physical models.
This invariance is the basic reason that a subspace comparison should not penalize a
pure retained-coordinate rotation. It also explains why projector agreement is weaker
than entrywise agreement of compressed operators.

\begin{theorem}[\texttt{BG-ALIGNMENT-THM-005}: Covariance of operator compression]
\label{thm:bg-alignment-005}
For every ambient operator $H$, frame $V$, and retained-space unitary $G$,
\[
  C_{V_G}(H)=G^\dagger C_V(H)G.
\]
Neither Hermiticity of $H$ nor orthonormality of $V$ is required.
\end{theorem}

\begin{proof}
Expanding the changed frame and reversing the conjugate transpose of its product gives
\[
  C_{V_G}(H)
  =(VG)^\dagger H(VG)
  =G^\dagger V^\dagger HVG
  =G^\dagger C_V(H)G.
\]
\end{proof}

A compressed operator is not invariant as a matrix: it transforms by unitary
similarity. Thus its individual entries are coordinate-dependent, while similarity-
invariant information such as its characteristic polynomial is not. This covariance is
why operator comparisons must first place both reduced matrices in a common retained
coordinate system. Selection of the aligning unitary $G$ is a separate problem.

\begin{theorem}[\texttt{BG-ALIGNMENT-THM-006}: Hermiticity of compression]
\label{thm:bg-alignment-006}
If $H^\dagger=H$, then for every frame $V$,
\[
  C_V(H)^\dagger=C_V(H).
\]
Frame orthonormality is not required for this algebraic closure.
\end{theorem}

\begin{proof}
Conjugate-transpose reversal and the Hermiticity of $H$ yield
\[
  C_V(H)^\dagger
  =(V^\dagger HV)^\dagger
  =V^\dagger H^\dagger V
  =V^\dagger HV
  =C_V(H).
\]
\end{proof}

When $H$ is a finite Hamiltonian, this result ensures that its compressed matrix remains
Hermitian and can therefore be treated as a valid finite Hermitian eigenproblem. The
closure is algebraic. Spectral approximation and subspace adequacy require additional
hypotheses.

\begin{theorem}[\texttt{BG-ALIGNMENT-THM-007}: Reconstruction is a projector sandwich]
\label{thm:bg-alignment-007}
For every frame $V$ and ambient operator $H$,
\[
  R_V(H)=P_VHP_V.
\]
Here $P_V$ need not be idempotent unless the frame is orthonormal.
\end{theorem}

\begin{proof}
This is exact reassociation of the matrix products:
\[
  R_V(H)
  =V(V^\dagger HV)V^\dagger
  =(VV^\dagger)H(VV^\dagger)
  =P_VHP_V.
\]
\end{proof}

The identity $R_V(H)=P_VHP_V$ exposes the ambient operator actually retained by the
finite reduction. It shows that compression followed by reconstruction is a two-sided
restriction to the chosen frame operator, rather than an unexplained smaller matrix.
For an orthonormal frame this is the familiar projected operator $P_VHP_V$; the theorem
itself carefully avoids assuming that $P_V$ is already idempotent.

\begin{theorem}[\texttt{BG-ALIGNMENT-THM-008}: Reconstruction invariance]
\label{thm:bg-alignment-008}
For every frame $V$, retained-space unitary $G$, and ambient operator $H$,
\[
  R_{V_G}(H)=R_V(H).
\]
\end{theorem}

\begin{proof}
Apply Theorem~\ref{thm:bg-alignment-007} to the changed frame and then use
Theorem~\ref{thm:bg-alignment-004}:
\[
  R_{V_G}(H)=P_{V_G}HP_{V_G}=P_VHP_V=R_V(H).
\]
Thus the retained matrix changes by unitary similarity, but its reconstructed ambient
operator does not change.
\end{proof}

Although the compressed matrix changes under $V\mapsto VG$, the ambient reconstructed
operator does not. This is the precise sense in which retained-coordinate gauge changes
alter representation but not the reconstructed reduced operator. It provides a
basis-independent comparison object and guards against treating a similarity transform
as a physical discrepancy.

\begin{theorem}[\texttt{BG-ALIGNMENT-THM-009}: The projector fixes its frame]
\label{thm:bg-alignment-009}
If $V^\dagger V=I_M$, then
\[
  P_VV=V.
\]
\end{theorem}

\begin{proof}
By associativity and the exact frame hypothesis,
\[
  P_VV=(VV^\dagger)V=V(V^\dagger V)=VI_M=V.
\]
\end{proof}

The projector acts as the identity on every retained column. This confirms that the
ambient projection leaves the represented retained states unchanged, a basic
consistency condition for interpreting those columns as the range of $P_V$. On the orthogonal
complement, the projector acts as zero rather than as the ambient identity.

\begin{theorem}[\texttt{BG-ALIGNMENT-THM-010}: Compression of the retained projector]
\label{thm:bg-alignment-010}
If $V^\dagger V=I_M$, then the compression of $P_V$ is the retained identity:
\[
  C_V(P_V)=I_M.
\]
\end{theorem}

\begin{proof}
Expanding $P_V$ and reassociating gives
\[
  C_V(P_V)
  =V^\dagger(VV^\dagger)V
  =(V^\dagger V)(V^\dagger V)
  =I_MI_M=I_M.
\]
\end{proof}

In retained coordinates, the projector becomes exactly $I_M$. This supplies a useful
normalization and regression identity: a correct compression of the retained projector
cannot introduce a nontrivial retained-space operator. It closes the relationship
between the ambient projector and its retained coordinate representation.

\begin{remark}[Scope of the theorem sequence]
These ten theorems concern exact finite matrices and equal-rank right-unitary coordinate
changes. Runtime frame construction, unequal-rank comparison, spectral adequacy, and
momentum-space regularity enter later layers with additional hypotheses.
\end{remark}
